% Generated draft: requires mathematical and source review.
\documentclass[11pt]{article}
\usepackage[margin=25mm]{geometry}
\usepackage{fontspec}
\setmainfont{Latin Modern Roman}
\usepackage{amsmath,amssymb,mathtools,unicode-math}
\setmathfont{Latin Modern Math}
\usepackage{enumitem,needspace,xurl,hyperref}
\setlength{\parindent}{0pt}
\setlength{\parskip}{5pt}
\setlength{\emergencystretch}{3em}
\newcommand{\N}{\mathbb{N}}
\newcommand{\Z}{\mathbb{Z}}
\newcommand{\Q}{\mathbb{Q}}
\newcommand{\R}{\mathbb{R}}
\newcommand{\C}{\mathbb{C}}
\newcommand{\F}{\mathbb{F}}
\newcommand{\A}{\mathbb{A}}
\newcommand{\PP}{\mathbb{P}}
\newcommand{\E}{\mathbb{E}}
\newcommand{\eps}{\varepsilon}
\newcommand{\dd}{\,\mathrm{d}}
\begin{document}
\section*{UnsolvedMath ID: 11000222}
\textbf{Source problem code:} catalogue entry
\subsection{Definitions and mathematical statement}
For a translation-surface model \((X,\omega)\), the usual Teichmüller flow is induced by \(g_t=\operatorname{diag}(e^t,e^{-t})\) on translation charts. The associated Gauss--Manin/Kontsevich--Zorich cocycle \(A_t\) acts on \(H^1(X;\R)\). For a specified invariant ergodic measure and integrable norm data, its Lyapunov exponents are the asymptotic logarithmic singular-value growth rates
\[
 \lambda_i=\lim_{t\to\infty}\frac1t\log s_i(A_t)
\]
at almost every point. The source does not explicitly give the measure, area/time normalization or whether the flow-derivative or cohomological cocycle spectrum is intended; its subsequent \(H^1\) discussion points to the latter.

Problem 15 asks to study the \emph{individual} exponents:
\begin{enumerate}
\item for all known \(\mathrm{SL}_2(\R)\)-invariant subvarieties;
\item for strata;
\item for strata in large genera as \(g\to\infty\).
\end{enumerate}
It also asks whether these exponents are related to characteristic numbers of natural bundles over suitable compactifications of strata. No bundles, compactifications, characteristic-number formula, genus sequence or meaning of ``all known'' is supplied. The motivation is Kontsevich's [Ko] differential-geometric formula for the sum of the first \(g\) exponents; the formula itself is not included, and the request must not be reduced to that sum.

The subsequent context invokes Calabi [Cb]: a real closed \(1\)-form \(\omega_0\) on a smooth real surface \(S\), with isolated zero set \(\Sigma\) and unspecified natural conditions, can be the real part of a holomorphic \(1\)-form \(\omega\) for some complex structure. The missing conditions are essential; the claim is not made here for arbitrary closed forms. The resulting \((X,\omega)\) need not be unique.

For generic resulting pairs, the source attributes to Oseledets theory [BaPe] transverse contracting and expanding Lagrangian subspaces
\[
 H^1(X;\R)=L_0\oplus L_1.
\]
A full contracting/expanding decomposition needs absence of zero exponents in the intended setting, not just the abstract Oseledets theorem; genericity, measure and this hypothesis are omitted. The source further reports that \(L_0\) depends only on \(\omega_0\), not its holomorphic realization, and is unchanged by small deformations preserving the indicated relative class. The printed class is \([\omega]\in H^1(S,\Sigma;\R)\), despite \(\omega\) being complex-valued; whether this means \([\omega_0]\), the real part, or another convention must be checked. Identifications of cohomology and the zero set during deformation are also missing. These reports define the intended implicit dynamical/topological object only after those qualifications are supplied.
\subsection{Short English statement}
Study individual Lyapunov exponents on the specified invariant loci, on strata and in large genus, and seek characteristic-number interpretations beyond the known sum formula. The source's measure, normalization and later contracting-Lagrangian hypotheses are missing and remain explicit source dependencies.
\subsection{Sources}
\begin{itemize}
\item Local UnsolvedMath registry, permanent ID \(11000222\), exact source code \texttt{catalogue entry}, Problem 15, ``Lyapunov exponents,'' PDF page 261. No individual proposer or proposed year is supplied.
\item Farb, edited volume \emph{Problems on Mapping Class Groups and Related Topics}: \url{https://www.math.uchicago.edu/~farb/papers/mcgbook.pdf}. The volume has 329 explicitly typeset numbered items across chapters, repeated local numbering distinguished by PDF page. The extract includes chapter/page furniture on billiards, flat and translation surfaces, page 255.
\item Internal references: Kontsevich [Ko] for the first-\(g\)-exponent sum formula; Calabi [Cb] for realization of closed forms under conditions; Oseledets theory via [BaPe]. Full bibliography entries and the precise formulas/hypotheses are absent.
\item Extraction pdf-font-structure, accessed 2026-07-29. Source evidence NEEDS\_REVIEW presents an open problem and calls for a current primary-literature audit. Supplied OPEN-TRIAGE summary names Kontsevich--Zorich and Forni, without complete citations, and says full computation remains open.
\item Dataset attribution: Ulam AI / UnsolvedMath Contributors, CC BY 4.0. Source rights NEEDS\_REVIEW: publicly accessible author/publisher copy, redistribution terms requiring review.
\end{itemize}
\end{document}
